\chapter{Masked autoencoding: validity, targets, remove-and-restore, erosion}\label{ch:masking}

\begin{onething}
Masked training is four small index computations done in order: decide which tokens hold real data
(average-pool the pixel mask, keep a token when more than half its pixels are valid), hide a fixed
fraction of the valid ones by sorting random numbers, physically delete the hidden tokens with a
gather and put them back as zeros with a scatter, and grade the model only on hidden, valid pixels
that are not too close to invalid ones.
\end{onething}

\Cref{ch:recon} established the rule that makes reconstruction useful for anomaly detection: a pixel
must never help reconstruct itself. The encoder (\cref{ch:encoder}) showed \emph{where} hiding happens:
at stage 1, between the patch embedding and the blocks. This chapter is the machinery behind that one
line of the encoder, and the masks that the trainer builds around it. All of it lives in one file,
\code{app/foundation_models/components/token_masking.py}, a class of static methods with no learnable
parameters, plus a few lines of the trainer,
\code{app/foundation_models/trainers/hyperspectral_segformer_mae_trainer.py}.

\section{The problem, and why the obvious fixes fail}\label{sec:masking:why}

You want the model to learn ``what normally goes here'' by predicting parts of a patch it cannot see.
Three obvious approaches each fail in a specific way.

\begin{enumerate}
\item \textbf{Hide pixels by setting them to zero.} A zeroed pixel still enters the patch embedding and
still becomes part of a token that attention can read. Worse, the patch already contains genuine
zeros: invalid pixels (cloud, no-data, swath edge) were zeroed before normalisation
(\cref{ch:bandpipeline}). The model could not tell ``hidden'' from ``missing''.
\item \textbf{Hide random pixels, not tokens.} A token is built from a 4 $\times$ 4 block. If you hide
half of a block's pixels, the token still contains the other half, and the model can interpolate from
inside the very token it is graded on.
\item \textbf{Grade the model everywhere.} Then it is graded on invalid pixels, where the ``truth'' is
a fill value, and on pixels it could see, where copying is a perfect answer.
\end{enumerate}

The design that survives is: hide whole tokens, choose them only among tokens that hold real data,
delete them from the encoder's input so nothing can attend to them, and grade only the pixels that
were both hidden and valid.

\begin{table}[htbp]\centering
\caption{The four steps of masking and the masks each produces. $B$ is the batch size, $N = 1024$
tokens for a 128 $\times$ 128 patch.}\label{tab:masking:map}
\small
\begin{tabular}{@{}clll@{}}\toprule
Step & Question & Output & Section \\\midrule
1 & Which tokens hold real data? & token mask $v$, \shape{(B, N)} & \ref{sec:masking:validity} \\
2 & Of those, which are hidden? & keep mask $k$, prediction mask $q$, \shape{(B, N)} & \ref{sec:masking:choose} \\
3 & How are they deleted and put back? & gather indices \shape{(B, N\_kept)} & \ref{sec:masking:gather} \\
4 & Which pixels are graded? & loss mask \shape{(B, 1, H, W)} & \ref{sec:masking:targets}, \ref{sec:masking:erosion} \\\bottomrule
\end{tabular}
\end{table}

\begin{analogy}
Think of a cloze test, the reading exercise in which words are blanked out of a paragraph and you must
fill them in. A good examiner blanks whole words, not half a word (a token, not a pixel). She never
blanks a word that was already illegible on the photocopy (an invalid token), because there is no
right answer to mark. She removes the blanked words from the sheet entirely rather than writing them
in faint ink (deletion, not zeroing). And she does not mark a filled-in word when the smudge from the
illegible word next to it covers half of it (erosion). Each of those rules is one section of this
chapter.
\end{analogy}

\section{Step 1: which tokens hold real data}\label{sec:masking:validity}

\subsection{The mismatch}

Validity arrives per \emph{pixel}: a tensor $m$ of shape \shape{(B, 1, H, W)} with
$m_{y,x} = 1$ where the pixel is a real measurement and $0$ where it is not. Masking works per
\emph{token}: one decision for each of the $N$ tokens. One stage-1 token covers a 4 $\times$ 4 block of
16 pixels (\cref{ch:encoder}, $r_1 = j_1 = 4$). When a block is partly valid, somebody has to make a
ruling.

\subsection{The rule}

\begin{definition}[Token validity]
Let $P_t$ be the set of 16 pixels of token $t$. The \index{valid fraction}\emph{valid fraction} of $t$
is
\begin{equation}\label{eq:masking:frac}
f_t = \frac{1}{\abs{P_t}} \sum_{(y,x) \in P_t} m_{y,x} = \frac{1}{16} \sum_{(y,x) \in P_t} m_{y,x} .
\end{equation}
Token $t$ is \index{token validity}\emph{valid} when $f_t > \tfrac12$, strictly:
\begin{equation}\label{eq:masking:valid}
v_t = \begin{cases} 1 & \text{if } f_t > 0.5,\\ 0 & \text{otherwise.}\end{cases}
\end{equation}
\end{definition}

Why is \cref{eq:masking:frac} a fraction? Because $m$ contains only 0s and 1s, the sum counts the
valid pixels, and dividing by 16 turns the count into a proportion. Sixteen valid pixels give $16/16 =
1$; eight give $8/16 = 0.5$; none gives $0$.

The code computes all $f_t$ at once with an \index{average pooling}\emph{average pool}, an operation
that slides a window over a grid and outputs the mean of each window
(\file{token_masking.py}, lines 112--127):
\begin{lstlisting}
token_fractions = F.avg_pool2d(pixel_mask.float(),     # (B, 1, 128, 128)
                               kernel_size=4, stride=4, padding=0)   # -> (B, 1, 32, 32)
token_valid = (token_fractions > 0.5).float()          # strict: 0.5 itself is invalid
token_mask  = token_valid.squeeze(1).flatten(1)        # (B, 32, 32) -> (B, 1024)
\end{lstlisting}
\code{squeeze(1)} removes the size-1 channel axis; \code{flatten(1)} joins the two grid axes into one
list of 1024 in \index{row-major order}\emph{row-major order} (row 0 left to right, then row 1, and so
on), the same order the tokens themselves use (\cref{ch:shapes}).

\begin{worked}[label={we:masking:validity}]{Token validity on a 4 \texorpdfstring{$\times$}{x} 4 pixel grid with 2 \texorpdfstring{$\times$}{x} 2 tokens}
To keep the arithmetic small, shrink the token to 2 $\times$ 2 pixels (pooling kernel 2, stride 2; the
real model uses 4 and 4). The pixel mask is
\[
m = \begin{pmatrix} 1 & 1 & 1 & 0\\ 1 & 1 & 0 & 0\\ 1 & 0 & 1 & 1\\ 1 & 0 & 1 & 1 \end{pmatrix}.
\]
The 2 $\times$ 2 token grid has four tokens, numbered row-major: $t_0$ top-left, $t_1$ top-right,
$t_2$ bottom-left, $t_3$ bottom-right. Do each on paper.

Token $t_0$, rows 0--1, columns 0--1: pixels $1, 1, 1, 1$.
\begin{align*}
f_0 &= \tfrac14 (1 + 1 + 1 + 1) = \tfrac44 = 1.00, & 1.00 &> 0.5 \;\Rightarrow\; v_0 = 1 .
\end{align*}
Token $t_1$, rows 0--1, columns 2--3: pixels $1, 0, 0, 0$.
\begin{align*}
f_1 &= \tfrac14 (1 + 0 + 0 + 0) = \tfrac14 = 0.25, & 0.25 &\not> 0.5 \;\Rightarrow\; v_1 = 0 .
\end{align*}
Token $t_2$, rows 2--3, columns 0--1: pixels $1, 0, 1, 0$.
\begin{align*}
f_2 &= \tfrac14 (1 + 0 + 1 + 0) = \tfrac24 = 0.50, & 0.50 &\not> 0.5 \;\Rightarrow\; v_2 = 0 .
\end{align*}
Token $t_3$, rows 2--3, columns 2--3: pixels $1, 1, 1, 1$.
\begin{align*}
f_3 &= \tfrac44 = 1.00 \;\Rightarrow\; v_3 = 1 .
\end{align*}
Flattened row-major: $v = [1, 0, 0, 1]$. Token $t_2$ is exactly half valid and is rejected:
\textbf{ties go to ``invalid''}. The same computation with the real 4 $\times$ 4 tokens on the course's
8 $\times$ 8 grid is in \cref{fig:masking:validity}; it ends in $[1, 0, 1, 0]$, again with one token
rejected at exactly $8/16$.
\end{worked}

The tie rule is the conservative choice: a token built half from no-data is not something you want to
grade the model on.

\coursefig[0.95\linewidth]{iw-17-token-validity}{Token validity on an 8 $\times$ 8 pixel mask with the
real stage-1 geometry (kernel 4, stride 4). (a) The pixel mask split into four 4 $\times$ 4 blocks.
(b) Average pooling gives each block's valid fraction: 16/16, 4/16, 12/16, 8/16. (c) The strict
threshold $> 0.5$ gives the token mask $[1, 0, 1, 0]$; the exactly-half block is invalid.}{fig:masking:validity}

\subsection{Exactly the same geometry as the patch embedding}

The pooling uses kernel 4, stride 4, padding 0: the stage-1 patch embedding's numbers
(\cref{ch:patchembed}). That is a requirement, not a coincidence. With a different window, a token's
\emph{validity} would be judged from a different set of pixels than the token's \emph{value} was
computed from, and you could grade the model on a token built entirely from no-data. The trainer and
the inferencer both declare the constants at the top of their files (\code{STAGE1_KERNEL_SIZE = 4},
\code{STAGE1_STRIDE = 4}, \code{STAGE1_PADDING = 0}; trainer lines 36--38, inferencer
\file{app/foundation_models/inferencers/segformer_mae_inferencer.py} lines 41--43) and pass them at
every call site.

\begin{sourcenote}[: padding defaults]
\code{pixel_mask_to_token_mask} defaults \code{padding=None} to \code{kernel_size // 2}
(\file{token_masking.py}, lines 107--108), which for kernel 4 would be 2 and would \emph{not} match
stage 1. Every caller passes \code{padding=0} explicitly, so the default is never used for stage 1. If
you add a caller, pass the constant.
\end{sourcenote}

\subsection{What happens to invalid tokens}

This is the part people get wrong. An invalid token ($v_t = 0$) is
\begin{itemize}
\item \textbf{not} deleted from the encoder's input,
\item \textbf{not} eligible to be hidden,
\item \textbf{not} graded.
\end{itemize}
It stays in the sequence carrying its zero-filled, normalised values (an invalid pixel sits near
$-1.22$ after normalisation, \cref{ch:bandpipeline}), so the encoder is told where the holes in the data
are rather than being asked to invent them.

\begin{keypoint}
The token mask $v$ says ``is this real data?''; it is a fact about the sensor. The keep mask of the
next section says ``shall the encoder see it?''; it is a choice. An invalid token has $v = 0$ and keep
$= 1$.
\end{keypoint}

\section{Step 2: choosing which valid tokens to hide}\label{sec:masking:choose}

\subsection{How many: the hiding count and its rounding}

Let $n_{\text{valid}} = \sum_t v_t$ be the number of valid tokens and $\rho$ the
\index{mask ratio}\emph{mask ratio}, the fraction of valid tokens to hide. The number hidden is
\begin{equation}\label{eq:masking:count}
n_{\text{hide}} = \lfloor \rho \, n_{\text{valid}} \rfloor ,
\end{equation}
because the code converts with \code{(num_valid * mask_ratio).long()} (\file{token_masking.py}, line
186), and converting a positive float to an integer in PyTorch truncates toward zero, which for
positive numbers is the floor. The ratio is a fraction of the \emph{valid} tokens, not of all tokens.

The count is deterministic: for a given patch the number hidden is fixed; only \emph{which} tokens are
hidden is random. Because every valid token is equally likely to be any of the hidden ones (shown
below), each valid token is hidden with probability
\begin{equation}\label{eq:masking:effective}
\Pr[\text{token $t$ hidden}] = \frac{n_{\text{hide}}}{n_{\text{valid}}} = \frac{\lfloor \rho\, n_{\text{valid}} \rfloor}{n_{\text{valid}}} \le \rho ,
\end{equation}
the \index{mask ratio!effective}\emph{effective} mask ratio. The floor can only lower it.

\begin{worked}[label={we:masking:rounding}]{Six tokens at \texorpdfstring{$\rho = 0.65$}{rho = 0.65}: what the rounding does}
Six valid tokens, $\rho = 0.65$ (the training value). On paper:
\begin{align*}
\rho\, n_{\text{valid}} &= 0.65 \times 6 && \text{multiply}\\
&= 3.9 ,\\
n_{\text{hide}} &= \lfloor 3.9 \rfloor = 3 && \text{truncate, do not round to 4}\\
n_{\text{keep}} &= 6 - 3 = 3 .
\end{align*}
The effective ratio is $3/6 = 0.50$, well below $0.65$. Rounding to nearest would have given 4 and an
effective $4/6 \approx 0.67$. For the real patch the effect is tiny:
\begin{align*}
0.65 \times 1024 &= 665.6, & \lfloor 665.6 \rfloor &= 665, & \frac{665}{1024} &\approx 0.6494 ,
\end{align*}
so a fully valid patch hides 665 tokens and keeps $1024 - 665 = 359$. The lesson: on tiny examples the
floor dominates; on real patches it costs a fraction of a token.
\end{worked}

Two edge cases follow from \cref{eq:masking:count}. A patch with one valid token hides
$\lfloor 0.65 \rfloor = 0$ tokens, and so does a patch with none: the loss mask is then empty and the
trainer returns a zero loss with zero counted pixels, which the training loop skips
(\file{hyperspectral_segformer_mae_trainer.py}, lines 249--250). In practice the trainer has already
dropped every patch with less than 40\% valid pixels (\code{MIN_VALID_PIXEL_FRACTION = 0.4}, line 34).

\begin{drift}
$\rho = 0.65$ in v0.2.0 (\file{configs/hyperspectral_segformer_exp_2.json}) and in every spectral-stem
run up to FT04\tagM{research/experiments/indradhanu-finetune/FT04-final-joint/config.json}; v0.1.0
used 0.50 (\file{configs/hyperspectral_segformer_exp_1.json}). The config class default is 0.5
(\file{app/models/training/training_config.py}, line 97), which no shipped run uses. As of 2026-10-29.
\end{drift}

\subsection{Which ones: a random permutation by sorting noise}

The obvious way to choose is a Python loop: for each patch, list the valid tokens, shuffle, take the
first $n_{\text{hide}}$. On a GPU processing 128 patches at once, with a different list length per
patch, that loop would stall the device on every batch. So the choice is written with array
operations only, identical work for every patch. That is why it looks indirect.

\begin{definition}[Permutation, argsort]
A \index{permutation}\emph{permutation} of $\{0, \dots, N-1\}$ is a reordering: a list $\pi$ in which
each index appears exactly once. The \index{argsort}\emph{argsort} of a list of numbers $a$ is the
permutation that sorts it: $\pi = \operatorname{argsort}(a)$ means $a_{\pi_0} \le a_{\pi_1} \le \dots
\le a_{\pi_{N-1}}$. Read $\pi_k$ as ``the index of the $k$-th smallest value''.
\end{definition}

The recipe (\file{token_masking.py}, lines 169--223):
\begin{lstlisting}
noise = torch.rand(B, N)                        # uniform in [0, 1), one per token
noise = noise + (1 - token_mask) * 2            # invalid tokens move to [2, 3)
num_to_mask = (token_mask.sum(1) * mask_ratio).long()          # (B,)
sorted_indices = noise.argsort(dim=1)           # (B, N): smallest noise first
positions = torch.arange(N).float().expand(B, -1)              # 0, 1, ..., N-1
pred_mask_sorted = (positions < num_to_mask.unsqueeze(1)).float()   # first n_hide places
pred_mask = torch.zeros(B, N)
pred_mask.scatter_(1, sorted_indices, pred_mask_sorted)        # back to token order
keep_mask = 1.0 - pred_mask
\end{lstlisting}

In words: give every token a random priority; push every invalid token's priority up by 2; sort by
priority; flag the first $n_{\text{hide}}$ places in the queue; carry the flags back to the tokens they
belong to.

\paragraph{Why every hidden set is equally likely.} The noise values are independent draws from the
same continuous distribution. Any particular ordering of $N$ such values is exactly as likely as any
other (relabelling the tokens does not change the joint distribution), and ties have probability zero.
So $\operatorname{argsort}(\text{noise})$ is a \emph{uniformly random permutation}: each of the $N!$
orders has probability $1/N!$. Taking its first $n_{\text{hide}}$ entries therefore picks a uniformly
random subset of size $n_{\text{hide}}$, which is what \cref{eq:masking:effective} assumed.

\paragraph{Why the $+2$ protects invalid tokens.} Valid tokens keep noise in $[0, 1)$; invalid tokens
get noise in $[2, 3)$. Every valid value is below every invalid one, so all valid tokens sort ahead of
all invalid tokens. Since $n_{\text{hide}} \le n_{\text{valid}}$, the first $n_{\text{hide}}$ places
are always valid tokens. Any shift of at least 1 would work; 2 is a clear margin. The valid tokens
among themselves are still in uniformly random order, because adding the same constant to every invalid
token does not touch their relative order.

\begin{worked}[label={we:masking:perm}]{Hiding three of six tokens by hand}
Six tokens, all valid, $\rho = 0.65$, so $n_{\text{hide}} = 3$ (\cref{we:masking:rounding}). Suppose
the noise is
\[
\text{noise} = [0.42,\ 0.91,\ 0.07,\ 0.66,\ 0.18,\ 0.55] \quad (\text{tokens } t_0, \dots, t_5) .
\]
No token is invalid, so the $+2$ adds nothing. Sort the values ascending and write under each value the
token it came from:
\[
\begin{array}{l|cccccc}
\text{sorted value} & 0.07 & 0.18 & 0.42 & 0.55 & 0.66 & 0.91\\
\text{came from token} & 2 & 4 & 0 & 5 & 3 & 1
\end{array}
\]
So $\pi = \operatorname{argsort}(\text{noise}) = [2, 4, 0, 5, 3, 1]$. Flag the first three places:
\[
\text{pred\_mask\_sorted} = [\,0 < 3,\ 1 < 3,\ 2 < 3,\ 3 < 3,\ 4 < 3,\ 5 < 3\,] = [1, 1, 1, 0, 0, 0] .
\]
Carry the flags back. \code{scatter_} writes \code{pred_mask[pi[k]] = pred_mask_sorted[k]}:
\begin{align*}
k = 0:&\quad \pi_0 = 2, & \text{pred\_mask}[2] &= 1\\
k = 1:&\quad \pi_1 = 4, & \text{pred\_mask}[4] &= 1\\
k = 2:&\quad \pi_2 = 0, & \text{pred\_mask}[0] &= 1\\
k = 3:&\quad \pi_3 = 5, & \text{pred\_mask}[5] &= 0\\
k = 4:&\quad \pi_4 = 3, & \text{pred\_mask}[3] &= 0\\
k = 5:&\quad \pi_5 = 1, & \text{pred\_mask}[1] &= 0
\end{align*}
Result: $q = \text{pred\_mask} = [1, 0, 1, 0, 1, 0]$ and $k = \text{keep\_mask} = 1 - q = [0, 1, 0, 1, 0, 1]$.
Tokens 0, 2 and 4, the three with the smallest noise, are hidden.
\end{worked}

\subsection{The inverse permutation}

The original masked autoencoder \citep{he2022masked} writes the same idea with two index lists:
\code{ids_shuffle}, the argsort of the noise (our $\pi$), and \code{ids_restore}, its
\index{permutation!inverse}\emph{inverse}. This code base never builds \code{ids_restore}; the
\code{scatter_} does its job. Seeing why is worth the effort, because it is the same fact that makes
remove-and-restore work in \cref{sec:masking:gather}.

\begin{definition}[Inverse permutation]
The inverse of a permutation $\pi$ is the permutation $\pi^{-1}$ with $\pi^{-1}_{\pi_k} = k$ for every
$k$. If $\pi$ answers ``which token sits in place $k$?'', then $\pi^{-1}$ answers ``in which place does
token $t$ sit?''. That is the \emph{rank} of token $t$.
\end{definition}

\begin{worked}[label={we:masking:inverse}]{Building ids\_restore for six tokens}
Take $\pi = [2, 4, 0, 5, 3, 1]$ from \cref{we:masking:perm}. Apply the definition one place at a time:
place $k$ holds token $\pi_k$, so token $\pi_k$ has rank $k$.
\begin{align*}
\pi_0 = 2 &\;\Rightarrow\; \pi^{-1}_2 = 0, & \pi_1 = 4 &\;\Rightarrow\; \pi^{-1}_4 = 1, & \pi_2 = 0 &\;\Rightarrow\; \pi^{-1}_0 = 2,\\
\pi_3 = 5 &\;\Rightarrow\; \pi^{-1}_5 = 3, & \pi_4 = 3 &\;\Rightarrow\; \pi^{-1}_3 = 4, & \pi_5 = 1 &\;\Rightarrow\; \pi^{-1}_1 = 5.
\end{align*}
Read off in token order: $\pi^{-1} = [2, 5, 0, 4, 1, 3]$. Check: token 2 had the smallest noise
(0.07) and has rank 0; token 1 had the largest (0.91) and has rank 5.

There is a second way to get it: $\pi^{-1} = \operatorname{argsort}(\pi)$. Sort $\pi = [2,4,0,5,3,1]$
ascending: the value 0 is at place 2, value 1 at place 5, value 2 at place 0, value 3 at place 4, value
4 at place 1, value 5 at place 3. So $\operatorname{argsort}(\pi) = [2, 5, 0, 4, 1, 3]$, the same
list. (The MAE code computes \code{ids_restore = argsort(ids_shuffle)} for this reason.)

Now unshuffle the flags with a \emph{gather} instead of a scatter:
\begin{align*}
q_t &= \text{pred\_mask\_sorted}\big[\pi^{-1}_t\big] && \text{look up token $t$'s place}\\
q &= [\,s_2,\ s_5,\ s_0,\ s_4,\ s_1,\ s_3\,] && \text{with } s = [1,1,1,0,0,0]\\
  &= [\,1,\ 0,\ 1,\ 0,\ 1,\ 0\,] .
\end{align*}
The same $q$ as the scatter produced. Scattering through $\pi$ and gathering through $\pi^{-1}$ are the
same operation.
\end{worked}

Why are they the same? A scatter through $\pi$ writes $q_{\pi_k} = s_k$ for every place $k$. Substitute
$t = \pi_k$, so $k = \pi^{-1}_t$: then $q_t = s_{\pi^{-1}_t}$ for every token $t$, which is the gather
through $\pi^{-1}$. One line of algebra, and it is the reason this code needs no \code{ids_restore}.

\begin{worked}[label={we:masking:plus2}]{The course example: two invalid tokens and the \texorpdfstring{$+2$}{+2}}
This is the example in the source's own docstring (\file{token_masking.py}, lines 147--163). Six tokens
with $v = [1, 1, 0, 1, 1, 0]$ (tokens 2 and 5 invalid) and $\rho = 0.5$.

Noise: $[0.30, 0.70, 0.10, 0.50, 0.20, 0.80]$. Token 2, which is invalid, drew the smallest value; left
alone it would be hidden first.

Add 2 where $v = 0$:
\begin{align*}
&[0.30,\ 0.70,\ 0.10,\ 0.50,\ 0.20,\ 0.80]\\
+\;&[0,\ 0,\ 2,\ 0,\ 0,\ 2]\\
=\;&[0.30,\ 0.70,\ 2.10,\ 0.50,\ 0.20,\ 2.80] .
\end{align*}
Count: $n_{\text{valid}} = 1+1+0+1+1+0 = 4$, $n_{\text{hide}} = \lfloor 4 \times 0.5 \rfloor = 2$.
Sort: $0.20 (t_4) < 0.30 (t_0) < 0.50 (t_3) < 0.70 (t_1) < 2.10 (t_2) < 2.80 (t_5)$, so
$\pi = [4, 0, 3, 1, 2, 5]$. Flags $[1, 1, 0, 0, 0, 0]$ scatter to $q[4] = 1$, $q[0] = 1$, all others 0:
\[
q = [1, 0, 0, 0, 1, 0], \qquad k = [0, 1, 1, 1, 0, 1] .
\]
Two of the four valid tokens are hidden; neither invalid token was touched; both invalid tokens have
$k = 1$ and stay in the sequence.
\end{worked}

\coursefig[0.95\linewidth]{iw-18-targets}{The worked example of \cref{we:masking:plus2} as one
picture. Invalid tokens (hatched) get $+2$ so they sort to the back. After an ascending sort the first
$\lfloor 4 \times 0.5 \rfloor = 2$ places (amber) are flagged and scattered back: tokens 0 and 4 are
hidden, and the keep mask is the exact complement of the prediction mask.}{fig:masking:targets}

There is no third state: every token is either hidden-and-graded ($q = 1$, $k = 0$) or
shown-and-not-graded ($q = 0$, $k = 1$), and $k = 1 - q$ always.

\subsection{Choosing the ratio}

The ratio sets the difficulty of the exam. At $\rho = 0.1$ nine tokens in ten are visible and filling
the tenth is interpolation; the model learns smoothing and, at inference, reconstructs anomalies well
from their surroundings, so they stop standing out. At $\rho = 0.9$ there is too little evidence, the
best strategy is a blurry average, and every pixel has a large residual. The masked-autoencoder
literature on photographs settled near 0.75 \citep{he2022masked}; this project uses 0.65.

\begin{sourcenote}[: inference does not use this function]
At inference every pixel needs a score, so random hiding is replaced by an exact complementary split
(a checkerboard, or a random half and its complement; \cref{ch:inference}). The invalid-token guarantee
is kept by multiplication instead of the $+2$: \code{pred_mask = token_validity * (1.0 - checker)}
(\file{segformer_mae_inferencer.py}). Only valid tokens can become targets either way.
\end{sourcenote}

\section{Step 3: gather and scatter, remove and restore}\label{sec:masking:gather}

\subsection{The two index operations}

\begin{definition}[Gather and scatter]
Given data $x$ of length $N$ and a list of indices $g = (g_0, \dots, g_{M-1})$:
\begin{itemize}
\item a \index{gather}\emph{gather} reads from the named positions: $y_k = x_{g_k}$ for $k = 0, \dots,
M-1$. The output is as long as the index list.
\item a \index{scatter}\emph{scatter} writes into the named positions: starting from a given $z$
(here all zeros), set $z_{g_k} = y_k$. Positions not named keep their starting value.
\end{itemize}
For tokens, each $x_t$ is a vector of $C$ numbers and the whole vector moves.
\end{definition}

A tiny example. With $x = [a, b, c, d, e]$ and $g = [3, 0, 4]$, the gather gives $[d, a, e]$. Scattering
$[d, a, e]$ through the same $g$ into $[0,0,0,0,0]$ writes $d$ at 3, $a$ at 0, $e$ at 4: $[a, 0, 0, d, e]$.
Gather then scatter with the same indices puts every gathered value back where it came from and leaves
zeros elsewhere.

\paragraph{The same thing as matrices.} Build the \index{selection matrix}\emph{selection matrix}
$\mat{S}$ of size $M \times N$ whose row $k$ is all zeros except a 1 in column $g_k$. Then
\begin{align*}
\text{gather:}\quad \vect{y} &= \mat{S}\vect{x}, &
\text{scatter into zeros:}\quad \vect{z} &= \mat{S}\T \vect{y} .
\end{align*}
Check the first on the tiny example: row 0 of $\mat{S}$ has its 1 in column 3, so
$y_0 = 0 \cdot a + 0 \cdot b + 0 \cdot c + 1 \cdot d + 0 \cdot e = d$. Two products tell the whole story:
\begin{itemize}
\item $\mat{S}\mat{S}\T = \mat{I}_M$: entry $(k, l)$ is the dot product of rows $k$ and $l$ of
$\mat{S}$, which is 1 when they point at the same column ($k = l$, since the $g_k$ are distinct) and 0
otherwise. Scatter then gather returns exactly what you scattered.
\item $\mat{S}\T\mat{S} = \diag(\text{keep})$: an $N \times N$ diagonal matrix with 1 at the gathered
positions and 0 elsewhere. Gather then scatter keeps the gathered tokens and zeroes the rest.
\end{itemize}
The second line \emph{is} remove-then-restore, with the encoder's blocks in between.

\begin{worked}[label={we:masking:smatrix}]{The selection matrix for three kept tokens of six}
From \cref{we:masking:perm}, $k = [0, 1, 0, 1, 0, 1]$: tokens 1, 3, 5 are kept. With $g = [1, 3, 5]$,
\[
\mat{S} = \begin{pmatrix} 0&1&0&0&0&0\\ 0&0&0&1&0&0\\ 0&0&0&0&0&1 \end{pmatrix}.
\]
Multiply $\mat{S}\T\mat{S}$ entry by entry: entry $(i, j)$ sums $S_{ki} S_{kj}$ over the three rows $k$.
It is nonzero only when column $i$ and column $j$ both hold a 1 in the same row, which happens only for
$i = j \in \{1, 3, 5\}$. So $\mat{S}\T\mat{S} = \diag(0, 1, 0, 1, 0, 1) = \diag(k)$.
\end{worked}

\subsection{remove\_tokens, line by line}

\begin{lstlisting}
num_kept = keep_mask.sum(dim=1).long()                       # (B,)
max_kept = num_kept.max().item()                             # one number for the batch
sorted_indices = keep_mask.argsort(dim=1, descending=True)   # (B, N): 1s first
gather_indices = sorted_indices[:, :max_kept]                # (B, max_kept)
idx = gather_indices.unsqueeze(-1).expand(-1, -1, C)         # (B, max_kept, C)
kept_tokens = torch.gather(tokens, dim=1, index=idx)         # (B, max_kept, C)
\end{lstlisting}
Sorting the keep mask descending puts every 1 before every 0, so the first \code{max_kept} indices are
the kept positions. \code{torch.gather} needs an index tensor of the output's shape, so each index is
repeated $C$ times by \code{expand} (a view, no copy): for $g = [1, 3, 5]$ and $C = 2$ the expanded index
is $[[1,1], [3,3], [5,5]]$, which pulls whole token vectors.

\begin{worked}[label={we:masking:rr}]{Remove and restore on six tokens}
Continue with $k = [0, 1, 0, 1, 0, 1]$, tokens $t_0, \dots, t_5$, each of $C = 2$ numbers.
\begin{align*}
\text{num\_kept} &= 0 + 1 + 0 + 1 + 0 + 1 = 3, \qquad \text{max\_kept} = 3,\\
\text{argsort descending of } k &= [1, 3, 5, 0, 2, 4] \quad (\text{the three 1s, then the three 0s}),\\
g &= [1, 3, 5],\\
\text{kept} &= [t_1, t_3, t_5], \quad \text{shape } (1, 3, 2) .
\end{align*}
The blocks turn these into $t_1', t_3', t_5'$. Restore starts from zeros of shape $(1, 6, 2)$ and
scatters through $g$:
\begin{align*}
\text{position } 1 &\leftarrow t_1', & \text{position } 3 &\leftarrow t_3', & \text{position } 5 &\leftarrow t_5',
\end{align*}
giving $[\,\vect{0},\ t_1',\ \vect{0},\ t_3',\ \vect{0},\ t_5'\,]$. The hidden positions 0, 2, 4 are
exactly zero vectors.
\end{worked}

\begin{sourcenote}[: the order of the kept tokens is not guaranteed, and does not matter]
The course says the descending sort keeps ties ``in index order''. \code{torch.argsort} is not stable by
default (\code{stable=False}), so among equal keys the order is implementation-defined. On this
machine (CPU, PyTorch 2.5.1) the small cases happen to come out in index order, $[1, 3, 5, 0, 2, 4]$
for the keep mask above (\file{research/textbook/tools/compute_masking_output.txt}), but on a GPU and
for $N = 1024$ nothing promises it. It does not
matter, for two reasons. First, restore scatters through the \emph{same} $g$, so every token returns to
its own slot whatever order it travelled in. Second, nothing in the stage-1 blocks depends on order when
tokens have been removed: the linear layers and LayerNorms act on each token separately, attention
takes its full-attention fallback, which treats the keys as a set, and Mix-FFN skips its depthwise
convolution (\cref{ch:encoder}). There is no positional encoding to disturb. In the language of
\cref{ch:attention}, a masked stage 1 is \emph{permutation-equivariant}: shuffling the input tokens
shuffles the outputs the same way.
\end{sourcenote}

\subsection{The documented wrinkle: max\_kept}

\code{max_kept} is the largest kept count in the batch. If patch A keeps 400 tokens and patch B keeps
380, both gather 400 indices, and B's last 20 come from the tail of its sort, which is its
\emph{hidden} tokens. Those 20 slip back into the encoder. The source says so: ``shorter items get a few
extra tokens from the sort tail (minor leakage)'' (\file{token_masking.py}, lines 301--304).

\begin{worked}[label={we:masking:leak}]{Leakage in a batch of two}
Patch A is fully valid: $n_{\text{valid}} = 1024$, $n_{\text{hide}} = 665$, kept 359. Patch B has
$n_{\text{valid}} = 1000$ and 24 invalid tokens:
\begin{align*}
n_{\text{hide}}^{B} &= \lfloor 0.65 \times 1000 \rfloor = \lfloor 650 \rfloor = 650,\\
\text{kept}^{B} &= 1024 - 650 = 374 \quad (\text{350 visible valid plus 24 invalid}).
\end{align*}
Then $\text{max\_kept} = \max(359, 374) = 374$, and patch A gathers $374 - 359 = 15$ indices from the
tail of its sort: 15 of its 665 hidden tokens, about $15/665 \approx 2.3\%$, are seen by the encoder
and still graded. Notice the direction: the patch with \emph{fewer} invalid tokens keeps fewer tokens
(its hiding count is larger) and is the one that leaks.
\end{worked}

The source calls it tolerable because counts differ only when validity differs, and the 40\%-valid
filter keeps batches fairly homogeneous. It is a real, small leak: on mixed batches with very different
cloud cover the fully valid patches leak most.

\subsection{restore\_tokens, and why zeros}

\begin{lstlisting}
full_tokens = torch.zeros(B, N, C)                                 # (B, 1024, 32)
scatter_indices = gather_indices.unsqueeze(-1).expand(-1, -1, C)   # (B, N_kept, C)
full_tokens.scatter_(1, scatter_indices, kept_tokens)              # kept slots written
\end{lstlisting}
Positions not written stay zero (\file{token_masking.py}, lines 402--417). The original masked
autoencoder fills them with a learned \emph{mask token} \citep{he2022masked}; this model does not, and
the encoder chapter gives the reason (the decoder fuses four maps and reads the gaps from F2--F4). The
round trip for one patch at $\rho = 0.65$:
\[
\begin{aligned}
&\shape{(B, 1024, 32)} \xrightarrow{\text{remove}} \shape{(B, 359, 32)} \xrightarrow{\text{2 blocks, LN}} \shape{(B, 359, 32)}\\
&\quad\xrightarrow{\text{restore}} \shape{(B, 1024, 32)} \xrightarrow{\text{reshape}} \shape{(B, 32, 32, 32)} .
\end{aligned}
\]

\coursefig[0.95\linewidth]{iw-19-remove-restore}{Hidden tokens are removed, not zeroed. The gather
keeps only the visible tokens (here 4 of 6), so the stage-1 blocks never see the hidden ones; after the
blocks the scatter writes the encoded tokens back into a zero tensor at their original positions.}{fig:masking:rr}

\paragraph{Why delete rather than zero?} Two reasons. \emph{Zero leakage}: a zeroed token is still a
key; after the pre-norm LayerNorm and the projections' biases it is a perfectly ordinary key and value,
and it takes a share of every softmax. In the unmasked path it would also enter the depthwise
convolution and the reduction convolution of its neighbours. A deleted token cannot influence anything. \emph{Cost}: the linear layers of stage 1 run on
359 tokens instead of 1024. (The attention scores do not get cheaper, because a masked stage 1 falls
back to full attention: \cref{ch:encoder}.)

\section{Step 4a: prediction targets and the loss mask}\label{sec:masking:targets}

The model returns a full reconstruction $\xhat$ of shape \shape{(B, 165, H, W)}; the trainer decides
which pixels count. It builds three pixel-level masks (\file{hyperspectral_segformer_mae_trainer.py},
lines 186--238):

\begin{enumerate}
\item The \textbf{pixel prediction mask} $Q$: each token's flag $q_t$ copied to its 16 pixels by
nearest-neighbour upsampling (\code{F.interpolate(..., mode="nearest")} from 32 $\times$ 32 to
128 $\times$ 128). Pixel $(y, x)$ belongs to token $(\lfloor y/4 \rfloor, \lfloor x/4 \rfloor)$, so
$Q_{y,x} = q_{\lfloor y/4 \rfloor, \lfloor x/4 \rfloor}$.
\item The \textbf{eroded validity mask} $E = \operatorname{erode}(m)$ of \cref{sec:masking:erosion}.
In training the erosion kernel is 1, which leaves $E = m$.
\item The \textbf{loss mask}, their product:
\begin{equation}\label{eq:masking:lossmask}
L_{y,x} = Q_{y,x} \cdot E_{y,x} .
\end{equation}
\end{enumerate}
A pixel is graded when it was hidden \emph{and} is valid (and, with erosion on, far enough from an
invalid pixel). The L1 part of the loss (\cref{ch:losses}) averages the per-pixel band-mean error over
the set $\Omega = \{(y,x) : L_{y,x} = 1\}$:
\begin{equation}\label{eq:masking:l1}
\mathcal{L}_{1} = \frac{1}{\abs{\Omega}} \sum_{(y,x) \in \Omega} \frac{1}{165} \sum_{c=1}^{165} \abs{\hat{x}_{c,y,x} - x_{c,y,x}} .
\end{equation}

\begin{worked}[label={we:masking:lossmask}]{The loss mask on the 4 \texorpdfstring{$\times$}{x} 4 example}
Return to \cref{we:masking:validity} with 2 $\times$ 2 tokens, $v = [1, 0, 0, 1]$. Two tokens are
valid; say $\rho = 0.65$ gives $\lfloor 0.65 \times 2 \rfloor = \lfloor 1.3 \rfloor = 1$ hidden token,
and the noise picks $t_3$ (bottom-right). So $q = [0, 0, 0, 1]$. Upsample to pixels, then multiply by
$m$ (erosion off):
\[
Q = \begin{pmatrix} 0&0&0&0\\0&0&0&0\\0&0&1&1\\0&0&1&1 \end{pmatrix},\quad
m = \begin{pmatrix} 1&1&1&0\\1&1&0&0\\1&0&1&1\\1&0&1&1 \end{pmatrix},\quad
L = Q \odot m = \begin{pmatrix} 0&0&0&0\\0&0&0&0\\0&0&1&1\\0&0&1&1 \end{pmatrix}.
\]
($\odot$ is the entry-by-entry product.) $\abs{\Omega} = 4$. Now suppose instead the course's 8
$\times$ 8 grid of \cref{fig:masking:validity} hid its token 2, which has 12 valid pixels of 16. Then
$Q$ has 16 ones in that block but $L$ has only $12$: \textbf{a hidden token is graded only on its valid
pixels}. A token can be valid (more than half) and still contain invalid pixels, and those are never
graded.
\end{worked}

\section{Step 4b: erosion of the mask border}\label{sec:masking:erosion}

\subsection{The problem}

Take a pixel that is perfectly valid but sits two pixels from a cloud edge. Its token's 4 $\times$ 4
block may include cloud pixels, which were zeroed and sit near $-1.22$ in normalised space. The
token's value is built partly from ground and partly from ``no data'', and the reconstruction there is
poor for a purely technical reason. In a scene this shows up as a bright ring around every cloud, every
water edge and every ragged swath edge in the anomaly map: the most common false-positive pattern in
this system. The fix is to shrink the valid region by a border strip and not count that strip in the
loss (training) or the score (inference). The pixels are still reconstructed; they do not count.

\subsection{Binary erosion and dilation}

\begin{definition}[Binary erosion and dilation]
Let $A$ be a set of pixels (here the valid ones) and $K_r$ the square of offsets
$\{(\delta_y, \delta_x) : \abs{\delta_y} \le r,\ \abs{\delta_x} \le r\}$, of side $k = 2r + 1$. The
\index{erosion}\emph{erosion} of $A$ keeps a pixel only if its whole square neighbourhood lies in $A$:
\[
A \ominus K_r = \{\, p : p + \delta \in A \text{ for every } \delta \in K_r \,\}.
\]
The \index{dilation}\emph{dilation} of a set $D$ adds every pixel within reach of it:
\[
D \oplus K_r = \{\, p : p + \delta \in D \text{ for some } \delta \in K_r \,\}.
\]
\end{definition}

Erosion and dilation are linked by complements. Write $A^c$ for the pixels not in $A$ (the invalid
ones). Then
\begin{equation}\label{eq:masking:duality}
(A \ominus K_r)^c = A^c \oplus K_r .
\end{equation}
Derive it in words, one step at a time:
\begin{align*}
p \notin A \ominus K_r
&\iff \text{not } \big(p + \delta \in A \text{ for every } \delta\big) && \text{definition of erosion}\\
&\iff p + \delta \notin A \text{ for some } \delta && \text{``not every'' is ``some not''}\\
&\iff p + \delta \in A^c \text{ for some } \delta && \text{definition of complement}\\
&\iff p \in A^c \oplus K_r && \text{definition of dilation.}
\end{align*}
So to erode the valid region you may \emph{grow} the invalid region and flip back. Growing a 0/1 mask is
a \index{max pooling}\emph{max pool} with stride 1: the window's maximum is 1 exactly when some pixel in
the window is 1. That is the whole implementation (\file{token_masking.py}, lines 64--74):
\begin{lstlisting}
padding = kernel_size // 2                     # radius r
invalid = 1.0 - mask                           # (B, 1, H, W): invalid -> 1
dilated_invalid = F.max_pool2d(invalid, kernel_size=kernel_size, stride=1, padding=padding)
eroded_mask = 1.0 - dilated_invalid            # flip back
\end{lstlisting}
The erosion radius is $r = \lfloor k/2 \rfloor$: kernel 5 removes 2 pixels, kernel 15 removes 7, and
kernel 1 removes nothing, because a 1 $\times$ 1 max pool returns its input. That is how ``off'' is
expressed without a special case.

\begin{worked}[label={we:masking:erode1d}]{Erosion in one dimension, kernel 5}
The mask $[0, 0, 1, 1, 1, 1, 1, 1, 1, 1, 0, 0]$ (indices 0--11), kernel 5, radius $\lfloor 5/2 \rfloor
= 2$. Flip: invalid $= [1, 1, 0, 0, 0, 0, 0, 0, 0, 0, 1, 1]$. Take the maximum over $[i - 2, i + 2]$,
clipped at the ends:
\begin{align*}
i = 0, 1, 2, 3&: \text{window contains index 0 or 1, max } 1\\
i = 4&: \text{indices } 2..6, \text{ all } 0, \text{ max } 0\\
i = 5, 6, 7&: \text{indices } 3..9 \text{ at most, all } 0, \text{ max } 0\\
i = 8, 9, 10, 11&: \text{window contains index 10 or 11, max } 1
\end{align*}
Dilated invalid: $[1,1,1,1,0,0,0,0,1,1,1,1]$. Flip back: $[0,0,0,0,1,1,1,1,0,0,0,0]$. The valid run was
indices 2--9 (8 pixels) and is now 4--7 (4 pixels): two removed from each end, as the radius promised.
\end{worked}

\begin{worked}[label={we:masking:erode2d}]{Erosion in two dimensions, kernel 3}
A 5 $\times$ 5 mask whose top row is invalid (a swath edge) and whose centre pixel is invalid (one bad
pixel). Kernel 3, radius 1: a pixel survives only if all 8 neighbours and itself are valid.
\[
m = \begin{pmatrix}0&0&0&0&0\\1&1&1&1&1\\1&1&0&1&1\\1&1&1&1&1\\1&1&1&1&1\end{pmatrix}
\quad\Longrightarrow\quad
\operatorname{erode}(m) = \begin{pmatrix}0&0&0&0&0\\0&0&0&0&0\\1&0&0&0&1\\1&0&0&0&1\\1&1&1&1&1\end{pmatrix}.
\]
Check three pixels by hand. Pixel $(1, 0)$: its window covers rows 0--2, and row 0 is invalid, so it
is eroded. Pixel $(2, 0)$: its window covers rows 1--3 and columns 0--1 (column $-1$ is outside the
image); none of those is invalid, so it survives. Pixel $(3, 1)$: its window covers rows 2--4 and
columns 0--2, which includes the bad centre $(2, 2)$, so it is eroded. The single bad pixel grew a
3 $\times$ 3 hole; the invalid row grew a second invalid row beneath it.
\end{worked}

\begin{sourcenote}[: the patch edge itself is not eroded]
\code{F.max_pool2d} pads with $-\infty$, not with zeros, so positions outside the image never count as
invalid. Erosion therefore grows only around \emph{invalid pixels}; the edge of a 128-pixel patch or of
the scene array is not trimmed unless the pixels there are invalid. Real swath and scene borders are
no-data pixels, so they are trimmed. The script \file{research/textbook/tools/compute_masking.py}
checks both cases: an all-valid 4 $\times$ 4 mask comes back unchanged with kernel 3.
\end{sourcenote}

\begin{sourcenote}[: the kernel must be odd]
Nothing validates oddness: the training field has \code{ge=1} and the inference field \code{gt=0}
(\file{app/models/training/training_config.py}, lines 101--106;
\file{app/models/training/inference_config.py}, lines 109--116). An even kernel breaks the shape. By the
output-size formula of \cref{ch:encoder} with stride 1 and padding $k/2$, a side $H$ becomes
$H + 2(k/2) - k + 1 = H + 1$, so the eroded mask would be one pixel larger than the image and the loss
mask product would fail. Both docstrings also say the kernel ``should match OPE kernel size (7 for
Stage 1)'' (\file{token_masking.py}, line 48; \file{inference_config.py}, line 115); stage 1's kernel
is 4 (\cref{ch:encoder}). Use odd kernels.
\end{sourcenote}

\subsection{Why throw the border strip away, and how wide}

A valid pixel within $r$ pixels of an invalid one shares its token, or a neighbouring token whose
patch-embedding window it feeds, with no-data. The stage-1 token is 4 pixels wide, so a radius of 3 is
enough to cover any token that touches an invalid pixel; larger radii also cover the wider windows of
stages 2--4 ($r_2 = 12$ pixels, \cref{ch:encoder}). Training uses radius 0 because training patches
come from scene interiors and are mostly fully valid: eroding every patch would throw away supervision
for little gain. Inference covers whole scenes with real cloud and swath edges, where a false ring costs
more than a few missing pixels, so it erodes by 7.

\begin{table}[htbp]\centering
\caption{The three erosions in the pipeline. A pixel must survive every one that applies to contribute
a score.}\label{tab:masking:erosions}
\begin{tabular}{@{}lllc@{}}\toprule
Erodes & Where & Field & Default kernel \\\midrule
sensor validity & training loss & \texttt{erosion\_kernel\_size} & 1 (off) \\
sensor validity & inference & \texttt{erosion\_kernel\_size} & 15 \\
analyst's region & scoring & \texttt{keep\_mask\_erosion\_kernel\_size} & 1 (off) \\\bottomrule
\end{tabular}
\end{table}

\begin{drift}
Defaults as of 2026-10-29: training erosion kernel 1\tagM{app/models/training/training_config.py},
inference 15\tagM{app/models/training/inference_config.py}; FT04 trained with 1
\tagM{research/experiments/indradhanu-finetune/FT04-final-joint/config.json}. The scoring erosion uses
SciPy's \code{binary_erosion} in \file{_anomaly_scoring_run.py} rather than a max pool.
\end{drift}

\coursefig[0.95\linewidth]{iw-20-erosion}{Erosion as a max pool, with the 1-D example of
\cref{we:masking:erode1d} (kernel 5). Flipping makes invalid pixels 1; a stride-1 max pool grows that
region by $\lfloor 5/2 \rfloor = 2$ pixels on each side; flipping back leaves the valid region shrunk by
2 pixels per edge.}{fig:masking:erosion}

\section{All four masks on one patch}\label{sec:masking:all}

\Cref{fig:masking:grids} puts the four steps on one small patch of 16 $\times$ 16 pixels with 4
$\times$ 4 tokens (16 tokens). A cloud covers rows 10--15, columns 11--15.

\begin{figure}[htbp]\centering
\begin{tikzpicture}[x=0.32cm,y=-0.32cm,font=\small]
  % token-level panel (a)
  \def\hid{{1,0,1,1,0,1,0,1,1,0,1,0,0,1,1,0}}
  \def\inv{{0,0,0,0,0,0,0,0,0,0,0,1,0,0,0,1}}
  \foreach \t in {0,...,15}{
    \pgfmathtruncatemacro{\r}{floor(\t/4)}\pgfmathtruncatemacro{\c}{mod(\t,4)}
    \pgfmathtruncatemacro{\h}{\hid[\t]}\pgfmathtruncatemacro{\iv}{\inv[\t]}
    \ifnum\iv=1
      \fill[maskgrey!45] (4*\c,4*\r) rectangle ++(4,4);
      \fill[pattern=north east lines,pattern color=maskgrey!90!black] (4*\c,4*\r) rectangle ++(4,4);
    \else\ifnum\h=1
      \fill[loss!55] (4*\c,4*\r) rectangle ++(4,4);
    \else
      \fill[spat!22] (4*\c,4*\r) rectangle ++(4,4);
    \fi\fi
    \node[font=\scriptsize] at (4*\c+2,4*\r+2) {\t};}
  \foreach \i in {0,4,...,16}{\draw[spat,thick] (\i,0) -- (\i,16); \draw[spat,thick] (0,\i) -- (16,\i);}
  \draw[densely dashed,thick] (11,16) -- (11,10) -- (16,10);
  \node[above] at (8,0) {(a) tokens};
  % pixel-level panel (b)
  \begin{scope}[shift={(20,0)}]
  \foreach \y in {0,...,15}{\foreach \x in {0,...,15}{
    \pgfmathtruncatemacro{\t}{4*floor(\y/4)+floor(\x/4)}
    \pgfmathtruncatemacro{\h}{\hid[\t]}
    \pgfmathtruncatemacro{\bad}{(\y>9)&&(\x>10)}
    \pgfmathtruncatemacro{\near}{(\y>8)&&(\x>9)}
    \ifnum\bad=1
      \fill[maskgrey!70] (\x,\y) rectangle ++(1,1);
    \else\ifnum\h=1
      \ifnum\near=1
        \fill[loss!25] (\x,\y) rectangle ++(1,1);
        \fill[pattern=crosshatch dots,pattern color=loss!90!black] (\x,\y) rectangle ++(1,1);
      \else
        \fill[loss!75] (\x,\y) rectangle ++(1,1);
      \fi
    \else
      \fill[white] (\x,\y) rectangle ++(1,1);
    \fi\fi}}
  \foreach \i in {0,...,16}{\draw[gray!50,line width=0.1pt] (\i,0) -- (\i,16); \draw[gray!50,line width=0.1pt] (0,\i) -- (16,\i);}
  \foreach \i in {0,4,...,16}{\draw[spat,thick] (\i,0) -- (\i,16); \draw[spat,thick] (0,\i) -- (16,\i);}
  \node[above] at (8,0) {(b) pixels: the loss mask};
  \end{scope}
  % legend
  \begin{scope}[shift={(0,17.5)},y=1cm,x=1cm]
    \fill[spat!22] (0,0) rectangle ++(0.35,-0.35); \node[right,font=\footnotesize] at (0.4,-0.17) {visible valid token};
    \fill[loss!55] (3.6,0) rectangle ++(0.35,-0.35); \node[right,font=\footnotesize] at (4.0,-0.17) {hidden token};
    \fill[maskgrey!45] (6.4,0) rectangle ++(0.35,-0.35);
    \fill[pattern=north east lines,pattern color=maskgrey!90!black] (6.4,0) rectangle ++(0.35,-0.35);
    \node[right,font=\footnotesize] at (6.8,-0.17) {invalid token};
    \fill[loss!75] (0,-0.6) rectangle ++(0.35,-0.35); \node[right,font=\footnotesize] at (0.4,-0.77) {graded pixel};
    \fill[loss!25] (3.6,-0.6) rectangle ++(0.35,-0.35);
    \fill[pattern=crosshatch dots,pattern color=loss!90!black] (3.6,-0.6) rectangle ++(0.35,-0.35);
    \node[right,font=\footnotesize] at (4.0,-0.77) {hidden, eroded};
    \fill[maskgrey!70] (6.4,-0.6) rectangle ++(0.35,-0.35); \node[right,font=\footnotesize] at (6.8,-0.77) {invalid pixel};
  \end{scope}
\end{tikzpicture}
\caption{The four masks on one 16 $\times$ 16 patch with 4 $\times$ 4 tokens, in the book's colours
(indigo spatial tokens, amber loss, grey invalid). A cloud (dashed) covers rows 10--15, columns 11--15.
(a) Tokens 11 and 15 are invalid (valid fractions $8/16$ and $0/16$); the other 14 are valid and
$\lfloor 0.65 \times 14 \rfloor = 9$ of them are hidden (an illustrative draw). (b) The loss mask with
erosion radius 1 (kernel 3): hidden pixels count only when valid and not next to the cloud. Worked in
\cref{we:masking:fourmasks}.}
\label{fig:masking:grids}
\end{figure}

\begin{worked}[label={we:masking:fourmasks}]{Counting the graded pixels in \cref{fig:masking:grids}}
Token $t$ covers rows $4\lfloor t/4 \rfloor$ to $+3$ and columns $4(t \bmod 4)$ to $+3$. The cloud is
rows $\ge 10$ and columns $\ge 11$.

\emph{Validity.} Token 10 (rows 8--11, columns 8--11) contains cloud at rows 10--11, column 11: 2
pixels, so $f_{10} = 14/16 = 0.875$, valid. Token 11 (rows 8--11, columns 12--15) has cloud in rows
10--11, all four columns: 8 pixels, so $f_{11} = 8/16 = 0.5$, \emph{invalid} by the strict rule. Token
14 (rows 12--15, columns 8--11) has cloud in column 11: 4 pixels, $f_{14} = 12/16 = 0.75$, valid.
Token 15 is all cloud, $f_{15} = 0$. So $n_{\text{valid}} = 16 - 2 = 14$.

\emph{Hiding.} $n_{\text{hide}} = \lfloor 0.65 \times 14 \rfloor = \lfloor 9.1 \rfloor = 9$; the draw
shown hides tokens $0, 2, 3, 5, 7, 8, 10, 13, 14$. That is $9 \times 16 = 144$ hidden pixels.

\emph{Without erosion}, the loss mask drops the hidden pixels that are invalid: 2 in token 10 and 4 in
token 14, so $144 - 2 - 4 = 138$ graded pixels.

\emph{With kernel 3} (radius 1), the invalid region grows by one pixel: rows $\ge 9$, columns $\ge 10$.
In token 10 that removes rows 9--11, columns 10--11: $3 \times 2 = 6$ pixels (the 2 invalid ones
included). In token 14, rows 12--15, columns 10--11: $4 \times 2 = 8$ pixels (the 4 invalid ones
included). Graded pixels: $144 - 6 - 8 = 130$. The 8 extra pixels lost to erosion are the border strip
next to the cloud.
\end{worked}

\begin{confusions}
\confusion{Is the token mask the same as the keep mask?}{No, and this is the most common mix-up. The
token mask says whether a token holds real data; the keep mask says whether the encoder is shown it.
Invalid tokens have token mask 0 and keep mask 1.}
\confusion{Is the mask ratio a fraction of all tokens?}{Of the valid ones: $\lfloor \rho\, n_{\text{valid}}
\rfloor$. A half-cloudy patch hides about half as many tokens in absolute terms.}
\confusion{Does 0.65 mean 65\% are hidden?}{Slightly fewer: the count is rounded down. A fully valid
patch hides 665 of 1024 (64.9\%); six valid tokens would hide 3 (50\%).}
\confusion{Why add exactly 2?}{Any shift of at least 1 puts invalid noise at or above 1, above every valid
value in $[0, 1)$. Two is a comfortable margin.}
\confusion{Where is ids\_restore?}{Nowhere. The scatter through the sort order is the same operation as a
gather through the inverse permutation (\cref{we:masking:inverse}), so it is never built.}
\confusion{Does a hidden token get graded on all 16 of its pixels?}{Only on the valid ones (and, with
erosion on, only those away from invalid pixels). A token can be valid with up to 7 invalid pixels.}
\confusion{Does the same mask get used every batch?}{No. \texttt{torch.rand} is called afresh each time, so
every batch of every epoch hides a different set.}
\confusion{Does erosion shrink the output image?}{No. The image keeps its size; eroded pixels are
reconstructed but not counted in the loss or score.}
\end{confusions}

\begin{checkyourself}
\item\label{ex:masking:1} A token's 4 $\times$ 4 block has 9 valid pixels. Compute $f_t$ and decide
whether it is valid. Then do the same for 8 valid pixels.
\item\label{ex:masking:2} A patch has 37 valid tokens and $\rho = 0.65$. How many are hidden, how many
tokens does the encoder see if the patch has 1024 tokens in all, and what is the effective mask ratio?
\item\label{ex:masking:3} Six tokens, $v = [1, 1, 1, 0, 1, 1]$, $\rho = 0.65$, noise $[0.9, 0.1,
0.05, 0.4, 0.6, 0.02]$. Work through the $+2$, the count, the argsort, the scatter and the keep mask.
Then compute the inverse permutation and check that gathering the sorted flags through it gives the
same prediction mask.
\item\label{ex:masking:4} \code{keep_mask = [1, 0, 1, 1]} for tokens $t_0..t_3$. Write the gather
indices, the kept tokens, the selection matrix $\mat{S}$, and the restored sequence after the blocks
produce $t_0', t_2', t_3'$.
\item\label{ex:masking:5} Erode $m = [0, 1, 1, 1, 1, 0]$ with kernel 3, by hand, using the
dilate-then-flip method. Then erode $[1, 1, 1, 1, 0, 1, 1, 1, 1]$ with kernel 5.
\item\label{ex:masking:6} In a batch of two, patch A has 1024 valid tokens and patch B 960. With
$\rho = 0.65$, compute each patch's kept count, \code{max_kept}, and how many hidden tokens leak into
which patch.
\end{checkyourself}

\nextchapter{ch:decoder}
